\chapter{The Model Zoo II: Recurrent Networks}
\label{ch:rnn}

The best single model of our campaign --- and of the 8th-place solution we
replicated --- is a recurrent network. Not because recurrence is
fashionable (it is not), but because three of its properties align exactly
with this genre's physics: causal-by-construction processing, cheap
statefulness over long days, and --- decisively --- graceful \emph{online
adaptation}. This chapter builds the architecture from its equations up to
the multi-branch variant that scored $+0.0092$, and collects the training
folklore that made it work.

\section{The data layout: a day is a batch}

Everything begins with a tensor convention. One trading day becomes a
tensor of shape
\[
(\underbrace{D}_{\text{symbols}},\ \underbrace{T}_{968 \text{ steps}},\
\underbrace{K}_{\text{features}}),
\]
symbols as the batch dimension, time as the sequence dimension. The
recurrent model reads each symbol's day left to right and emits a
prediction per step. Three consequences of this convention, all
load-bearing:

\begin{itemize}
\item \textbf{Causality is structural.} A left-to-right RNN cannot see
  $t+1$ at time $t$. The entire class of within-day leakage bugs
  (Chapter~\ref{ch:leakage}) is excluded by construction --- one reason to
  love recurrence at review time.
\item \textbf{The metric becomes the loss.} With one day per batch, the
  batch-level $1 - \Rw$ is a stable, day-normalized training loss
  (Chapter~\ref{ch:metric}).
\item \textbf{Hidden state = accumulated intraday context.} By afternoon,
  $h_t$ has digested the whole morning --- regime, volatility, the running
  history of every feature --- without any feature engineering.
\end{itemize}

\section{The cells}

For reference and for reading the variants below, the two standard gated
cells. GRU:
\begin{align*}
z_t &= \sigma(W_z x_t + U_z h_{t-1}) &\text{(update gate)}\\
r_t &= \sigma(W_r x_t + U_r h_{t-1}) &\text{(reset gate)}\\
\tilde h_t &= \tanh(W_h x_t + U_h (r_t \odot h_{t-1}))\\
h_t &= (1 - z_t) \odot h_{t-1} + z_t \odot \tilde h_t
\end{align*}
LSTM adds a separate cell state $c_t$ with input/forget/output gates
$i_t, f_t, o_t$:
\[
c_t = f_t \odot c_{t-1} + i_t \odot \tanh(W_c x_t + U_c h_{t-1}),
\qquad h_t = o_t \odot \tanh(c_t).
\]
Gates matter here for a concrete reason: they are \emph{learned forgetting
schedules}, and in a drifting, noisy environment the model's ability to
discard stale context is as important as its ability to accumulate it. A
consistent minor finding of our runs: \textbf{the LSTM edged the GRU in
every online-refit comparison} (e.g.\ $+0.0092$ vs $+0.0089$ at matched
budgets) despite the GRU's better static starts --- the extra forget
machinery earns its parameters exactly when the model must keep rewriting
its own state day after day. Minor, consistent, and the reason our
production bag leans LSTM.

\section{ModelR: the multi-branch architecture that won its tier}

The 8th-place architecture (``ModelR'' in our registry,
\code{gru\_modelr}/\code{lstm\_modelr}) is worth studying as a piece of
design economy. It is \emph{four} parallel recurrent stacks, each ending in
a linear head that predicts one \textbf{auxiliary responder}
(\S\ref{sec:auxtargets}); the four branch outputs are then combined by a
single linear layer into the final target prediction:
\[
\hat a^{(i)} = \text{RNN}_i(x_{1:t}) \in \mathbb R, \quad i = 1..4;
\qquad
\hat y_t = w^\top \big(\hat a^{(1)}, \dots, \hat a^{(4)}\big) + b .
\]
The training loss sums the weighted-$R^2$ losses of the target \emph{and}
each auxiliary head. Reading the design:

\begin{itemize}
\item Each branch is forced to become an expert on one related quantity
  (responders at other windows; synthetic forward extensions). The
  branches cannot collapse into one another --- they answer to different
  supervision.
\item The final prediction is a \emph{learned ensemble of four internal
  experts}. The variance-reduction logic of
  Chapter~\ref{ch:ensembling} operates \emph{inside} the model.
\item The auxiliary losses inject high-SNR gradient
  (\S\ref{sec:auxtargets}); our nowcast-head extension raises that SNR
  further by adding realized-window targets ($R^2 \approx 0.5$ tasks).
\end{itemize}

Measured: gru\_modelr $+0.0089$ online / lstm\_modelr $+0.0092$ online on
280 training days --- against $+0.0075$ for the same-size single-branch
GRU. The branches pay.

\section{Training folklore that mattered}

The published details of the reference solution plus our replication notes,
in the order they bit us:

\begin{itemize}
\item \textbf{Optimizer:} AdamW, lr $10^{-3}$, weight decay $10^{-2}$;
  gradient clipping at global norm 1.0 (non-negotiable with heavy-tailed
  targets --- a single wild day otherwise produces a gradient that
  scrambles the network).
\item \textbf{Batch = one date} (all symbols); days shuffled between
  epochs, order within a day never touched.
\item \textbf{Early stopping} on validation weighted-$R^2$ with patience
  $\sim$10 (Chapter~\ref{ch:diagnostics} for why patience must be long at
  this SNR); typical convergence 15--25 epochs.
\item \textbf{Size:} hidden sizes in the low hundreds. Our capacity
  sweep was flat: doubling the LSTM width moved the score by less than
  seed noise. \emph{Data and noise, not capacity, bind this problem} ---
  budget engineering hours accordingly.
\item \textbf{Depth over width, mildly:} 2--3 stacked recurrent layers
  with dropout $\sim$0.1 between them; deeper stacks gave nothing.
\item \textbf{Warm starts} for the bagging workflow: bag members
  (Chapter~\ref{ch:ensembling}) resume from a shared checkpoint and
  diverge via data subsets and seeds --- most of the cost of a fresh fit,
  saved.
\end{itemize}

\section{Online refitting: where the RNN earns its keep}

Chapter~\ref{ch:online} treats online learning fully; here, the
architectural point. Each morning, the competition hands you yesterday's
labels. The RNN update is one (or a few) gradient step(s) on yesterday's
day-batch:

\begin{lstlisting}
def update(self, X_day, y_day, w_day):        # yesterday, now labeled
    opt = AdamW(self.net.parameters(), lr=self.lr_refit)
    loss = weighted_r2_loss(self.net(X_day), y_day, w_day)
    loss.backward(); clip_grad_norm_(self.net.parameters(), 1.0)
    opt.step()
\end{lstlisting}

Measured effect on this family: \textbf{$+0.003$ to $+0.005$}, every
recurrent variant, every configuration --- the largest single lever in the
entire campaign, larger than any feature block and any architecture swap.
The same daily step barely moves attention models ($+0.0003$) and can
\emph{destroy} hybrid ones at the wrong learning rate
(Chapter~\ref{ch:online}'s cliff table). Why recurrence adapts so
gracefully is worth one sentence of speculation grounded in the gate
algebra: a gated cell's behavior is smoothly parameterized --- small weight
updates produce small, well-conditioned changes in its temporal filtering
--- whereas attention's routing (softmax over comparisons) can reorganize
discontinuously under the same step size. Whatever the mechanism, the
empirical rule is unambiguous: \emph{if the protocol delivers delayed
labels, recurrent models convert them into score at higher efficiency than
anything else we tested.}

\begin{jsbox}[: the numbers in one place]
280 training days, 100-day tail, identical harness:
xgb $+0.0068$ (static);
gru $+0.0075$ online;
gru\_modelr $+0.0089$ online;
lstm\_modelr $+0.0092$ online (refit lr $10^{-3}$; the sweep over refit
lr spans $+0.0063$ static $\to$ $+0.0092$ at $10^{-3}$ $\to$ $+0.0070$ at
$3\times10^{-3}$ --- the optimum is a plateau, not a point, but both
edges are real cliffs).
\end{jsbox}

\section{Failure modes and their signatures}

\begin{itemize}
\item \textbf{Hidden-state staleness across regime breaks}: a model
  trained on old data carries confident state into a new regime; static
  scores on recent tails go \emph{negative} (we measured $-0.002$ for
  800-day-old models). Remedy: recency windows + online refit, not bigger
  models.
\item \textbf{Exploding updates on wild days}: without clipping, one
  heavy-tailed day's refit step undoes weeks of fitting. Signature: online
  score cliff on a specific calendar date. Remedy: clip both offline and
  online.
\item \textbf{Silent shape bugs}: batching (symbols, time, features) makes
  it easy to transpose symbol and time and still run. Signature: the
  shuffle test of Chapter~\ref{ch:leakage} \emph{passes} (score survives
  within-day shuffling) when it should fail. Run it once per architecture.
\end{itemize}

\begin{drillbox}
Implement the smallest version of the big result: a one-layer GRU (hidden
32) on any drifting panel, trained on a 200-step window. Score its frozen
checkpoint on the next 50 steps; then score the same checkpoint with one
SGD step per revealed step at lr $\in \{0, 10^{-4}, 10^{-3}, 10^{-2}\}$.
Plot score versus refit lr. You will reproduce, in miniature, the
plateau-and-cliffs shape that Chapter~\ref{ch:online} elevates into a
design law.
\end{drillbox}
